% 作者题证的编辑转录副本，不是独立下载的上游原件。
% 来源：https://terrytao.wordpress.com/2011/09/27/254a-notes-3-haar-measure-and-the-peter-weyl-theorem/

\begin{theorem}[Existence and uniqueness of Haar measure; Theorem 3]
Let $G$ be a $\sigma$-compact locally compact Hausdorff group. There exists a
nonzero left-invariant Radon measure on $G$, unique up to multiplication by
a positive scalar. The analogous assertion holds for right-invariant
measures.
\end{theorem}
\paragraph{Measure convention and representation theorem.}
Here a Radon measure is a Borel measure finite on compact sets, inner
regular and outer regular. We use the Riesz representation theorem:
a positive linear functional on the real space $C_c(G;\R)$ is integration
against a unique Radon measure. Write $I_\mu(f)=\int_G f\,d\mu$ and
$\tau(y)f(x)=f(y^{-1}x)$.

\begin{proof}[Uniqueness]
Let $\mu,\nu$ be left Haar measures. It suffices to show
\[
 I_\nu(f)I_\mu(g)=I_\mu(f)I_\nu(g)\qquad(f,g\in C_c(G)).
\]
Fix $f,g$ and $\epsilon>0$. Uniform continuity on their compact supports
gives an identity neighbourhood $U_\epsilon$ such that
$|f(xy)-f(x)|\leq\epsilon$ and $|g(xy)-g(x)|\leq\epsilon$ for all $x\in G$,
$y\in U_\epsilon$. The neighbourhoods may be chosen inside one fixed compact
set. Every nonempty open set has positive Haar measure. Urysohn's lemma
therefore gives a nonnegative $\psi_\epsilon\in C_c(U_\epsilon)$ with
$\int\psi_\epsilon\,d\mu=1$. Thus
\[
 \int_G f(xy)\psi_\epsilon(y)\,d\mu(y)=f(x)+O(\epsilon)
\]
uniformly in $x$, with supports contained in a fixed compact set.
Integrating against $\nu$, then using left invariance and Fubini, gives
\[
\begin{aligned}
 I_\nu(f)+O(\epsilon)
 &=\int_G\int_G f(xy)\psi_\epsilon(y)\,d\mu(y)\,d\nu(x)\\
 &=\int_G\int_G f(y)\psi_\epsilon(x^{-1}y)\,d\mu(y)\,d\nu(x)\\
 &=\int_G f(y)\left(\int_G\psi_\epsilon(x^{-1})\,d\nu(x)\right)d\mu(y)\\
 &=I_\mu(f)\int_G\psi_\epsilon(x^{-1})\,d\nu(x).
\end{aligned}
\]
The same factor occurs for $g$. Eliminating it yields
$I_\nu(f)I_\mu(g)-I_\mu(f)I_\nu(g)=O(\epsilon)$.
Letting $\epsilon\to0$ proves the assertion, and Riesz representation
identifies $\nu$ as a scalar multiple of $\mu$. Nonzero positivity makes
that scalar positive.
\end{proof}

\begin{proof}[Existence]
It suffices to construct a nonnegative functional on $C_c(G)^+$ that is
positively homogeneous, additive and left-invariant, and is nonzero.
Fix a nonzero $f_0\in C_c(G)^+$. We first construct functionals
$I=I_{f_1,\ldots,f_n,\epsilon}$, for finite families $f_i\in C_c(G)^+$,
that satisfy homogeneity, left invariance, $I(f_0)=1$, a bound
$0\leq I(f)\leq K(f)$ independent of the chosen family and $\epsilon$, and
\[
 |I(f_i+f_j)-I(f_i)-I(f_j)|\leq\epsilon.
\]

Choose a small $\delta>0$ and an identity neighbourhood $U$ so that
$|f_i(xy)-f_i(x)|\leq\delta$ for all $x\in G$, $y\in U$, $0\leq i\leq n$.
Choose a nonzero $\psi\in C_c(G)^+$ supported in $U$. For $f\in C_c(G)^+$
define the covering number
\[
 [f:\psi]=\inf\left\{\sum_{k=1}^Kc_k:
 f\leq\sum_{k=1}^Kc_k\tau(y_k)\psi,\quad c_k\geq0,\ y_k\in G\right\}.
\]
Compactness of the support gives $[f:\psi]<\infty$, and
$[f_0:\psi]>0$ since $f_0\ne0$. Put
\[
 I(f)=\frac{[f:\psi]}{[f_0:\psi]}.
\]
This is homogeneous, left-invariant, monotone and normalized at $f_0$.
Furthermore
\[
 [f:\psi]\leq[f:f_0]\,[f_0:\psi],\qquad
 I(f+g)\leq I(f)+I(g).
\]
Hence the uniform bound holds with $K(f)=[f:f_0]$.

It remains to prove approximate additivity from below. Fix $i,j$. Choose a
cover
\[
 f_i(x)+f_j(x)\leq\sum_{k=1}^Kc_k\psi(y_k^{-1}x),\qquad
 \sum_kc_k\leq\left(I(f_i+f_j)+\frac\epsilon2\right)[f_0:\psi].
\]
Split $c_k=c_k'+c_k''$ by
\[
 c_k'=c_k\frac{f_i(y_k)+\delta}{f_i(y_k)+f_j(y_k)+2\delta},\qquad
 c_k''=c_k\frac{f_j(y_k)+\delta}{f_i(y_k)+f_j(y_k)+2\delta}.
\]
Whenever $\psi(y_k^{-1}x)\ne0$, uniform continuity gives
$|f_i(y_k)-f_i(x)|\leq\delta$ and $|f_j(y_k)-f_j(x)|\leq\delta$.
Consequently
\[
 \frac{f_i(y_k)+\delta}{f_i(y_k)+f_j(y_k)+2\delta}
 \geq\frac{f_i(x)}{f_i(x)+f_j(x)+4\delta}.
\]
Therefore
\[
 \sum_kc_k'\psi(y_k^{-1}x)+4\delta
 \geq\frac{f_i(x)(f_i(x)+f_j(x))}{f_i(x)+f_j(x)+4\delta}+4\delta
 \geq f_i(x).
\]
The analogous estimate holds for $f_j$ and the coefficients $c_k''$.
Take a fixed $g\in C_c(G)^+$ equal to $1$ on the supports of all the $f_i$.
The preceding estimates can be written with $4\delta g$ in place of the
constant $4\delta$: on the support this is the same bound, and outside the
support the required left-hand side is zero. Subadditivity then yields
\[
\begin{aligned}
 I(f_i)&\leq\frac{\sum_kc_k'}{[f_0:\psi]}+4\delta I(g),\\
 I(f_j)&\leq\frac{\sum_kc_k''}{[f_0:\psi]}+4\delta I(g).
\end{aligned}
\]
Thus
\[
 I(f_i)+I(f_j)\leq I(f_i+f_j)+\frac\epsilon2+8\delta I(g)
 \leq I(f_i+f_j)+\frac\epsilon2+8\delta K(g).
\]
Choosing $\delta$ small enough proves the desired approximate additivity
simultaneously for the finite family.

To pass to an exact functional, consider the compact product
\[
 \mathcal K=\prod_{f\in C_c(G)^+}[0,K(f)].
\]
For each finite family and $\epsilon>0$, impose the closed conditions of
homogeneity, left invariance, normalization and the specified approximate
additivity conditions. The functionals just constructed show that these
closed sets have the finite intersection property: for finitely many such
requirements take the union of the families and the smallest tolerance.
Tychonoff's theorem gives a point in their common intersection. This
functional is exactly additive, since the tolerances are arbitrarily small.
Extend it by differences to $C_c(G;\R)$. By the Riesz representation theorem
it is integration against a Radon measure $\mu$. Its normalization makes
$\mu$ nonzero, and the left invariance of the functional makes $\mu$
left-invariant. Inversion, $E\mapsto E^{-1}$, gives the right-invariant
version, completing the theorem.
\end{proof}
